\documentclass{article}
\usepackage[utf8]{inputenc}
\usepackage{amsmath}
\usepackage{amssymb}
\usepackage{amsfonts}

\begin{document}

\section*{Insiemi Numerici}

\begin{itemize}
    \item \textbf{Numeri Naturali:} $\mathbb{N} = \{0, 1, 2, 3, 4, \dots \}$ (Interi non negativi)
    \item \textbf{Numeri Interi:} $\mathbb{Z} = \{ \dots, -2, -1, 0, 1, 2, \dots \}$ (Interi positivi e negativi)
    \item \textbf{Numeri Razionali:} $\mathbb{Q} = \left\{ \frac{p}{q} \mid p, q \in \mathbb{Z}, q \neq 0 \right\}$ (Classi di equivalenza di frazioni, es. $\frac{1}{2} \equiv \frac{2}{4} \equiv \frac{3}{6}$)
    \item \textbf{Numeri Reali:} $\mathbb{R}$ (Include $\mathbb{Q}$ e molto altro, es. $\sqrt{2}$, $\pi$).
\end{itemize}

Catena di inclusioni: $\mathbb{N} \subset \mathbb{Z} \subset \mathbb{Q} \subset \mathbb{R}$.

\subsection*{Intervalli di $\mathbb{R}$}
Definizione: $I \subset \mathbb{R}$ si dice \textbf{intervallo} se dati $x, y \in I$ con $x < y$, risulta che per ogni $z$ tale che $x < z < y$, allora $z \in I$.
(In parole povere: in un intervallo non ci sono "buchi").

\textbf{Esempi:}
\begin{itemize}
    \item $A = \{ x \in \mathbb{R} : 2 < x < 4 \}$ è un intervallo.
    \item $B = \{ x \in \mathbb{R} : 1 \le x < 2 \text{ oppure } 6 < x < 7 \}$ \textbf{non} è un intervallo. (Controesempio: $x=1.5, y=6.5$, $z=4$ sta in mezzo ma $z \notin B$).
    \item $C = \mathbb{Z}$ non è un intervallo.
\end{itemize}

\textbf{Notazione:}
Dati $a, b \in \mathbb{R}$ con $a < b$:
\begin{itemize}
    \item $[a, b] = \{ x \in \mathbb{R} : a \le x \le b \}$ (Chiuso)
    \item $(a, b) = \{ x \in \mathbb{R} : a < x < b \}$ (Aperto)
    \item $[a, b) = \{ x \in \mathbb{R} : a \le x < b \}$
    \item $(a, b] = \{ x \in \mathbb{R} : a < x \le b \}$
    \item $[a, +\infty) = \{ x \in \mathbb{R} : x \ge a \}$ (Semiretta chiusa)
    \item $(a, +\infty) = \{ x \in \mathbb{R} : x > a \}$ (Semiretta aperta)
    \item $(-\infty, b] = \{ x \in \mathbb{R} : x \le b \}$
    \item $(-\infty, +\infty) = \mathbb{R}$
\end{itemize}

\section*{Funzioni}
Una funzione è una terna di oggetti $A, B, f$.
\begin{itemize}
    \item $A$: Dominio
    \item $B$: Codominio
    \item $f$: Legge che associa ad ogni elemento di $A$ \textbf{uno e un solo} elemento di $B$.
\end{itemize}
Notazione: $f: A \to B$.

\subsection*{Grafico e Immagine}
\begin{itemize}
    \item \textbf{Prodotto Cartesiano:} $A \times B = \{ (a,b) : a \in A, b \in B \}$.
    \item \textbf{Grafico:} $\text{graph}(f) = \{ (a,b) \in A \times B : b = f(a) \}$.
    \item \textbf{Immagine:} Dato $D \subseteq A$, $f(D) = \{ f(x) : x \in D \}$.
    \item L'immagine di tutta la funzione è $f(A) = \text{Imm}(f)$.
\end{itemize}

\subsection*{Proprietà delle Funzioni}
\begin{enumerate}
    \item \textbf{Iniettiva:} Se $x_1, x_2 \in A$ con $x_1 \neq x_2 \implies f(x_1) \neq f(x_2)$.
    \textit{Criterio grafico:} Ogni retta orizzontale incontra il grafico al massimo in un punto.
    \item \textbf{Suriettiva:} Se $\forall y \in B, \exists x \in A$ tale che $f(x) = y$. (Ovvero $\text{Imm}(f) = B$).
    \textit{Criterio grafico:} Ogni retta orizzontale incontra il grafico almeno in un punto.
    \item \textbf{Biettiva (o Biunivoca):} Se è sia iniettiva che suriettiva.
\end{enumerate}

\textbf{Funzione Inversa:}
Se $f: A \to B$ è biettiva, esiste $f^{-1}: B \to A$ tale che $f^{-1}(b) = a \iff f(a) = b$.
Esempio: $f(x) = x^2$ su $[0, +\infty)$ è biettiva. L'inversa è $f^{-1}(y) = \sqrt{y}$.

\end{document}
\documentclass{article}
\usepackage[utf8]{inputenc}
\usepackage{amsmath}
\usepackage{amssymb}
\usepackage{amsfonts}

\begin{document}

\section*{Insiemi Numerici}

\begin{itemize}
    \item \textbf{Numeri Naturali:} $\mathbb{N} = \{0, 1, 2, 3, 4, \dots \}$ (Interi non negativi)
    \item \textbf{Numeri Interi:} $\mathbb{Z} = \{ \dots, -2, -1, 0, 1, 2, \dots \}$ (Interi positivi e negativi)
    \item \textbf{Numeri Razionali:} $\mathbb{Q} = \left\{ \frac{p}{q} \mid p, q \in \mathbb{Z}, q \neq 0 \right\}$ (Classi di equivalenza di frazioni, es. $\frac{1}{2} \equiv \frac{2}{4} \equiv \frac{3}{6}$)
    \item \textbf{Numeri Reali:} $\mathbb{R}$ (Include $\mathbb{Q}$ e molto altro, es. $\sqrt{2}$, $\pi$).
\end{itemize}

Catena di inclusioni: $\mathbb{N} \subset \mathbb{Z} \subset \mathbb{Q} \subset \mathbb{R}$.

\subsection*{Intervalli di $\mathbb{R}$}
Definizione: $I \subset \mathbb{R}$ si dice \textbf{intervallo} se dati $x, y \in I$ con $x < y$, risulta che per ogni $z$ tale che $x < z < y$, allora $z \in I$.
(In parole povere: in un intervallo non ci sono "buchi").

\textbf{Esempi:}
\begin{itemize}
    \item $A = \{ x \in \mathbb{R} : 2 < x < 4 \}$ è un intervallo.
    \item $B = \{ x \in \mathbb{R} : 1 \le x < 2 \text{ oppure } 6 < x < 7 \}$ \textbf{non} è un intervallo. (Controesempio: $x=1.5, y=6.5$, $z=4$ sta in mezzo ma $z \notin B$).
    \item $C = \mathbb{Z}$ non è un intervallo.
\end{itemize}

\textbf{Notazione:}
Dati $a, b \in \mathbb{R}$ con $a < b$:
\begin{itemize}
    \item $[a, b] = \{ x \in \mathbb{R} : a \le x \le b \}$ (Chiuso)
    \item $(a, b) = \{ x \in \mathbb{R} : a < x < b \}$ (Aperto)
    \item $[a, b) = \{ x \in \mathbb{R} : a \le x < b \}$
    \item $(a, b] = \{ x \in \mathbb{R} : a < x \le b \}$
    \item $[a, +\infty) = \{ x \in \mathbb{R} : x \ge a \}$ (Semiretta chiusa)
    \item $(a, +\infty) = \{ x \in \mathbb{R} : x > a \}$ (Semiretta aperta)
    \item $(-\infty, b] = \{ x \in \mathbb{R} : x \le b \}$
    \item $(-\infty, +\infty) = \mathbb{R}$
\end{itemize}

\section*{Funzioni}
Una funzione è una terna di oggetti $A, B, f$.
\begin{itemize}
    \item $A$: Dominio
    \item $B$: Codominio
    \item $f$: Legge che associa ad ogni elemento di $A$ \textbf{uno e un solo} elemento di $B$.
\end{itemize}
Notazione: $f: A \to B$.

\subsection*{Grafico e Immagine}
\begin{itemize}
    \item \textbf{Prodotto Cartesiano:} $A \times B = \{ (a,b) : a \in A, b \in B \}$.
    \item \textbf{Grafico:} $\text{graph}(f) = \{ (a,b) \in A \times B : b = f(a) \}$.
    \item \textbf{Immagine:} Dato $D \subseteq A$, $f(D) = \{ f(x) : x \in D \}$.
    \item L'immagine di tutta la funzione è $f(A) = \text{Imm}(f)$.
\end{itemize}

\subsection*{Proprietà delle Funzioni}
\begin{enumerate}
    \item \textbf{Iniettiva:} Se $x_1, x_2 \in A$ con $x_1 \neq x_2 \implies f(x_1) \neq f(x_2)$.
    \textit{Criterio grafico:} Ogni retta orizzontale incontra il grafico al massimo in un punto.
    \item \textbf{Suriettiva:} Se $\forall y \in B, \exists x \in A$ tale che $f(x) = y$. (Ovvero $\text{Imm}(f) = B$).
    \textit{Criterio grafico:} Ogni retta orizzontale incontra il grafico almeno in un punto.
    \item \textbf{Biettiva (o Biunivoca):} Se è sia iniettiva che suriettiva.
\end{enumerate}

\textbf{Funzione Inversa:}
Se $f: A \to B$ è biettiva, esiste $f^{-1}: B \to A$ tale che $f^{-1}(b) = a \iff f(a) = b$.
Esempio: $f(x) = x^2$ su $[0, +\infty)$ è biettiva. L'inversa è $f^{-1}(y) = \sqrt{y}$.

\end{document}
